\chapter{Experimental Work}
\label{chap:experimental_work}

\section{Hardware Setup}

The experimental validation of the designed controllers was carried out using the WHEELTEC Z100 Mini-Bike platform described in Section~\ref{sec:minibike}.
The primary hardware components of the system are summarized as follows:

\begin{itemize}
	\item \textbf{Microcontroller:} STM32F103C8T6 (ARM Cortex-M3)
	\item \textbf{IMU Sensor:} MPU6050 (3-axis gyroscope and accelerometer)
	\item \textbf{Servo Motor:} WHEELTEC Z100 steering servo
	\item \textbf{Rear Wheel Motor:} DC motor with incremental encoder
	\item \textbf{Motor Driver:} AT8236 H-bridge circuit
	\item \textbf{Power Supply:} 2S lithium-ion battery (7--12 V)
	\item \textbf{Wireless Communication:} Bluetooth module (USART3)
\end{itemize}

Data acquisition and control commands were handled by the STM32 microcontroller operating at a sampling frequency of $100$~Hz.
Communication with the host computer for debugging and data logging was achieved through a UART serial interface.

\section{Software Implementation}

The control algorithms were implemented on the STM32 using C++ programming with HAL Drivers.
Key implementation steps included:
\begin{itemize}
	\item Reading IMU data (lean angle and angular velocity) via I2C communication.
	\item Processing encoder pulses to estimate forward velocity.
	\item Applying a complementary filter to estimate the lean angle more robustly.
	\item Implementing the state-feedback controller (LQR) and switching later to the Data-Driven MPC control law.
	\item Controlling the steering servo via PWM signal modulation.
	\item Logging real-time data through USART1 to a PC terminal for offline analysis.
	\item \todo{Bluetooth}
\end{itemize}

The real-time computation of the Data-Driven MPC control law involved solving a semi-definite programming (SDP) problem.
Due to computational limitations of the STM32, the SDP optimization was first performed offline, and the resulting state-feedback gain matrix was embedded into the embedded controller as a static gain.

\section{Data Collection}

Two types of data were collected during the experiments:
\begin{itemize}
	\item \textbf{Open-loop Data:} System response without feedback control to characterize the dynamics and validate the bicycle model.
	\item \textbf{Closed-loop Data:} System response under the LQR and Data-Driven MPC controllers.
\end{itemize}

The following signals were recorded:
\begin{itemize}
	\item Lean angle ($\phi$)
	\item Lean rate ($\dot{\phi}$)
	\item Steering angle ($\delta$)
	\item Control input (steering rate command)
	\item Motor encoder readings (optional)
\end{itemize}

Data was logged at a sampling rate of $100$~Hz, synchronized with control execution.

\section{Experimental Procedure}

The experimental procedure involved the following steps:
\begin{enumerate}
	\item Calibrate the IMU sensor to eliminate initial offsets.
	\item Start the bicycle with an initial forward velocity manually achieved by a push.
	\item Apply the selected controller (LQR or Data-Driven MPC).
	\item Record the system's lean angle and steering angle evolution during the stabilization process.
	\item Introduce manual disturbances by slightly perturbing the bicycle during motion.
	\item Stop the experiment after achieving stabilization or falling.
\end{enumerate}

Each experimental run was repeated multiple times to ensure reproducibility and to evaluate the robustness of the controller under different initial conditions and disturbances.

\section{Challenges Encountered}

Several practical challenges were encountered during the experimental phase:
\begin{itemize}
	\item \textbf{Sensor Noise:} The MPU6050 exhibited significant noise at low velocities, requiring careful filtering and calibration.
	\item \textbf{Timing and Delays:} Limited computational resources led to minor delays in control signal update, impacting the system's responsiveness.
	\item \textbf{Motor Saturation:} The servo motor exhibited saturation effects at high steering rates, limiting the maximum achievable control torque.
	\item \textbf{Balance Midpoint Drift:} Fine adjustments of the potentiometer on the main board were necessary to maintain correct balance points without modifying the control code.
\end{itemize}

Despite these challenges, stable experimental results were obtained, validating the effectiveness of the proposed data-driven control strategy.
